% =====================================================================
%  RAPPORT DE PROJET DE FIN D'ÉTUDES — FICHIER UNIQUE
%  Winaity Template Builder — France Téléphone / Groupe Winvest Capital
%
%  * Un seul fichier, aucun dossier, aucun \input.
%  * Compilation : pdfLaTeX uniquement (PAS de biber). Compile 2 fois
%    pour que la table des matières et les renvois \ref se mettent à jour.
%  * Sur Overleaf : Menu -> Compiler = pdfLaTeX.
% =====================================================================
\documentclass[12pt,a4paper,oneside]{report}

% ---- Encodage & langue ----
\usepackage[utf8]{inputenc}
\usepackage[T1]{fontenc}
\usepackage[french]{babel}
\usepackage{mathptmx}         % Times New Roman (texte + math), exigé par ESPRIT
\usepackage{textcomp}

% ---- Mise en page (conforme ESPRIT) ----
\usepackage[margin=2.5cm]{geometry}
\usepackage{setspace}
\setstretch{1.15}             % interligne 1.15
\usepackage{indentfirst}
\setlength{\parindent}{0.5cm} % alinéa de première ligne 0,5 cm
\setlength{\parskip}{0.4em}   % petit espace entre paragraphes

% ---- Figures & tableaux ----
\usepackage{graphicx}
\usepackage{float}
\usepackage{caption}
\usepackage{array}
\usepackage{tabularx}
\usepackage{booktabs}
\usepackage{multirow}
\usepackage{xcolor}
\usepackage{amssymb}          % \checkmark
\usepackage{pgfgantt}         % diagramme de Gantt (planning) — charge aussi TikZ
\usepackage{pdfpages}         % inclure la couverture/pages décoratives ESPRIT (PDF)
\usepackage{listings}         % code en annexe
\lstset{
  basicstyle=\ttfamily\footnotesize,
  breaklines=true,
  frame=single,
  columns=fullflexible,
  showstringspaces=false,
  keepspaces=true,
  xleftmargin=1em
}

% ---- Divers ----
\usepackage{enumitem}
\usepackage{hyperref}
\hypersetup{colorlinks=true,linkcolor=black,citecolor=black,urlcolor=blue}
\usepackage{fancyhdr}

% ---- En-têtes ----
\pagestyle{fancy}
\fancyhf{}
\fancyhead[L]{\leftmark}
\fancyfoot[R]{\thepage}
\renewcommand{\headrulewidth}{0.4pt}
\setlength{\headheight}{28pt}
\addtolength{\topmargin}{-16pt}

% ---- Placeholder de figure (à remplacer par \includegraphics) ----
\newcommand{\figplaceholder}[3]{%
  \begin{figure}[H]
    \centering
    \fbox{\parbox[c][#3][c]{0.7\linewidth}{\centering\itshape%
      [\,EMPLACEMENT FIGURE\,]\\ #1}}
    \caption{#1}
    \label{#2}
  \end{figure}%
}

% ---- Couleur ESPRIT ----
\definecolor{espritred}{RGB}{150,22,30}

% ---- Image si le fichier existe, sinon un cadre "à remplacer" ----
% Usage : \imgorbox{chemin.png}{largeur}{hauteurCadre}{texteCadre}
\newcommand{\imgorbox}[4]{%
  \IfFileExists{#1}{\includegraphics[width=#2]{#1}}%
  {\fbox{\parbox[c][#3][c]{#2}{\centering\itshape\small #4}}}%
}

\begin{document}

% =====================================================================
%  PAGE DE GARDE  (remplace par la page de garde officielle ESPRIT)
% =====================================================================
% ---------------------------------------------------------------------
%  PAGE DE GARDE (modèle ESPRIT)
%  -> Dépose les images dans un dossier images/ :
%       images/esprit-logo.png            (logo « esprit — se former autrement »)
%       images/logo-entreprise.png        (logo France Téléphone / Winvest)
%       images/esprit-accreditations.png  (bandeau UNESCO/CDIO/EUR-ACE/CGE)
%     Tant qu'une image manque, un cadre « à remplacer » s'affiche.
% ---------------------------------------------------------------------
% ===== MÉTHODE A (RECOMMANDÉE, rendu identique) =====================
% Remplis le modèle ESPRIT dans Word, exporte-le en PDF sous
%   images/page-de-garde.pdf   -> il est inséré ici tel quel (pixel-parfait).
% ===== MÉTHODE B (secours) : version LaTeX si le PDF n'est pas fourni =
\IfFileExists{images/page-de-garde.pdf}{%
  \includepdf[pages=-,fitpaper=true]{images/page-de-garde.pdf}%
}{%
\begin{titlepage}
\thispagestyle{empty}
% Fond décoratif ESPRIT optionnel : exporte la page du modèle SANS son
% texte (garde les triangles + bandeau) en images/esprit-fond.png
\IfFileExists{images/esprit-fond.png}{%
  \begin{tikzpicture}[remember picture,overlay]
    \node at (current page.center)
      {\includegraphics[width=\paperwidth,height=\paperheight]{images/esprit-fond.png}};
  \end{tikzpicture}}{}
\centering

\imgorbox{images/esprit-logo.png}{6cm}{2.2cm}{[ Logo ESPRIT ]}\\[1cm]

{\color{espritred}\Huge\bfseries 2025 -- 2026}\\[0.35cm]
{\Huge\bfseries PROJET DE FIN D'ÉTUDES}\\[1cm]

{\setlength{\fboxsep}{0pt}%
\noindent\colorbox{espritred}{\makebox[\textwidth][c]{%
  \rule[-0.4cm]{0pt}{1.2cm}\color{white}\LARGE\bfseries
  DIPLÔME NATIONAL D'INGÉNIEUR}}}\\[1cm]

{\Large\bfseries SPÉCIALITÉ : INFORMATIQUE}\\[1.2cm]

{\LARGE\bfseries Winaity Template Builder}\\[0.4cm]
{\large\itshape Conception et développement d'une plateforme SaaS
multi-entreprises de création de modèles de communication assistée par IA}\\[1.4cm]

\begin{flushleft}\large
\hspace{1.5cm}\textbf{Réalisé par :} [À COMPLÉTER — Prénom Nom]\\[0.45cm]
\hspace{1.5cm}\textbf{Encadré par :}\\[0.3cm]
\hspace{2.3cm}\textbf{Encadrant ESPRIT :} [À COMPLÉTER]\\[0.3cm]
\hspace{2.3cm}\textbf{Encadrant Entreprise :} [À COMPLÉTER]
\end{flushleft}

\vspace{0.5cm}
\hfill\imgorbox{images/logo-entreprise.png}{5cm}{2.4cm}{[ LOGO ENTREPRISE ]}\hspace{1.5cm}

\vfill
\imgorbox{images/esprit-accreditations.png}{\textwidth}{1.8cm}{[ Bandeau d'accréditations ESPRIT : UNESCO/UniTwin -- CDIO -- EUR-ACE/Cti -- Conférence des Grandes Écoles ]}
\end{titlepage}
}

% ===== PAGE DÉCORATIVE ESPRIT (toque de diplômé) ====================
% Exporte cette page du modèle en images/esprit-decorative.pdf (ou .png)
\IfFileExists{images/esprit-decorative.pdf}{%
  \includepdf[pages=-,fitpaper=true]{images/esprit-decorative.pdf}%
}{\IfFileExists{images/esprit-decorative.png}{%
  \thispagestyle{empty}%
  \begin{tikzpicture}[remember picture,overlay]
    \node at (current page.center)
      {\includegraphics[width=\paperwidth,height=\paperheight]{images/esprit-decorative.png}};
  \end{tikzpicture}\clearpage}{}}

% ---------------------------------------------------------------------
%  FORMULAIRE DE VALIDATION (à insérer APRÈS la page de garde)
%  Reconstruit intégralement en LaTeX — à imprimer, signer et scanner.
% ---------------------------------------------------------------------
\thispagestyle{empty}
\vspace*{1cm}
\noindent Je valide le dépôt du rapport PFE relatif à l'étudiant nommé
ci-dessous / \emph{I validate the submission of the student's report~:}

\vspace{0.4cm}
\noindent\hspace{0.8cm}\textbullet\hspace{0.3cm}Nom \& Prénom / \emph{Name \& Surname} : \dotfill

\vspace{0.8cm}
\noindent\fbox{\begin{minipage}{\dimexpr\linewidth-2\fboxsep-2\fboxrule\relax}
\vspace{0.2cm}
\centering{\large Encadrant Entreprise / \emph{Business site Supervisor}}
\vspace{0.5cm}
\begin{flushleft}
\hspace{0.5cm}\textbullet\hspace{0.3cm}Nom \& Prénom / \emph{Name \& Surname} : \dotfill
\end{flushleft}
\centering\textbf{Cachet \& Signature / \emph{Stamp \& Signature}}
\vspace{3cm}
\end{minipage}}

\vspace{0.8cm}
\noindent\fbox{\begin{minipage}{\dimexpr\linewidth-2\fboxsep-2\fboxrule\relax}
\vspace{0.2cm}
\centering{\large Encadrant Académique / \emph{Academic Supervisor}}
\vspace{0.5cm}
\begin{flushleft}
\hspace{0.5cm}\textbullet\hspace{0.3cm}Nom \& Prénom / \emph{Name \& Surname} : \dotfill
\end{flushleft}
\centering\textbf{Signature / \emph{Signature}}
\vspace{3cm}
\end{minipage}}

\vspace{0.8cm}
\noindent\small Ce formulaire doit être rempli, signé et \underline{scanné} /
\emph{This form must be completed, signed and \underline{scanned}.}\par
\noindent\small Ce formulaire doit être introduit après la page de garde /
\emph{This form must be inserted after the cover page.}
\clearpage

\pagenumbering{roman}

% =====================================================================
%  DÉDICACE
% =====================================================================
\chapter*{Dédicace}
\addcontentsline{toc}{chapter}{Dédicace}
\thispagestyle{empty}
\begin{center}
\itshape
\vspace{2cm}

Je dédie ce modeste travail

\vspace{1cm}

À mes chers parents,\\
pour leur amour inconditionnel, leurs sacrifices et leurs encouragements
constants qui m'ont accompagné tout au long de mon parcours.
Que ce travail soit le témoignage de ma profonde reconnaissance.

\vspace{1cm}

À toute ma famille et à mes amis,\\
pour leur soutien et leur présence dans les moments les plus exigeants.

\vspace{1cm}

À tous ceux qui, de près ou de loin,\\
ont contribué à faire de ce projet une réussite.

\vspace{2cm}

[À COMPLÉTER — Signature]
\end{center}

% =====================================================================
%  REMERCIEMENTS
% =====================================================================
\chapter*{Remerciements}
\addcontentsline{toc}{chapter}{Remerciements}
\thispagestyle{empty}

Au terme de ce projet de fin d'études, je tiens à exprimer ma gratitude
à toutes les personnes qui ont contribué, chacune à sa manière, à
l'aboutissement de ce travail.

J'adresse mes remerciements les plus sincères à mon encadrant
professionnel, \textbf{[À COMPLÉTER]}, au sein de \textbf{France Téléphone}
(groupe \textbf{Winvest Capital}), pour la confiance qu'il m'a accordée,
pour son accompagnement tout au long du stage et pour la richesse des
échanges techniques qui ont orienté les choix de conception de la
plateforme.

Je tiens également à remercier chaleureusement mon encadrant académique,
\textbf{[À COMPLÉTER]}, pour sa disponibilité, ses conseils avisés et sa
rigueur méthodologique, qui m'ont permis de structurer ma démarche et de
mener ce projet à son terme dans les meilleures conditions.

Mes remerciements s'étendent à l'ensemble de l'équipe de
\textbf{France Téléphone} pour leur accueil, leur esprit de collaboration
et l'environnement de travail stimulant dans lequel j'ai évolué.

J'exprime enfin ma reconnaissance aux membres du jury pour l'honneur qu'ils
me font en acceptant d'évaluer ce travail, ainsi qu'à l'ensemble du corps
enseignant de \textbf{ESPRIT} pour la qualité de la formation dispensée
tout au long de mon cursus.

% =====================================================================
%  SOMMAIRE + LISTES
% =====================================================================
\tableofcontents
\listoffigures
\listoftables

% =====================================================================
%  GLOSSAIRE
% =====================================================================
\chapter*{Glossaire}
\addcontentsline{toc}{chapter}{Glossaire}

\begin{description}[leftmargin=3.2cm,style=nextline,font=\bfseries]
  \item[API] Application Programming Interface (interface de programmation applicative)
  \item[CQRS] Command Query Responsibility Segregation (séparation commandes / requêtes)
  \item[CRM] Customer Relationship Management (gestion de la relation client)
  \item[CSS] Cascading Style Sheets
  \item[DnD] Drag and Drop (glisser-déposer)
  \item[DTO] Data Transfer Object (objet de transfert de données)
  \item[gRPC] gRPC Remote Procedure Call (framework d'appel de procédure distante)
  \item[HT] Hors Taxes
  \item[HTML] HyperText Markup Language
  \item[IA] Intelligence Artificielle
  \item[JSON] JavaScript Object Notation
  \item[JWT] JSON Web Token
  \item[LLM] Large Language Model (grand modèle de langage)
  \item[M2M] Machine to Machine (communication de service à service)
  \item[MCP] Model Context Protocol (protocole d'exposition d'outils à un modèle)
  \item[MJML] Mailjet Markup Language (langage de balisage pour e-mails responsives)
  \item[OCR] Optical Character Recognition (reconnaissance optique de caractères)
  \item[ORM] Object-Relational Mapping (mapping objet-relationnel)
  \item[PDF] Portable Document Format
  \item[RBAC] Role-Based Access Control (contrôle d'accès basé sur les rôles)
  \item[REST] Representational State Transfer
  \item[SaaS] Software as a Service (logiciel en tant que service)
  \item[SSO] Single Sign-On (authentification unique)
  \item[TTC] Toutes Taxes Comprises
  \item[TVA] Taxe sur la Valeur Ajoutée
  \item[UI/UX] User Interface / User Experience
  \item[vLLM] Moteur d'inférence haute performance pour grands modèles de langage
\end{description}

% =====================================================================
%  CORPS
% =====================================================================
\cleardoublepage
\pagenumbering{arabic}

% ---------------------------------------------------------------------
%  INTRODUCTION GÉNÉRALE
% ---------------------------------------------------------------------
\chapter*{Introduction générale}
\addcontentsline{toc}{chapter}{Introduction générale}
\markboth{Introduction générale}{Introduction générale}

La communication constitue aujourd'hui un pilier stratégique de toute
organisation. Qu'il s'agisse d'une campagne promotionnelle par e-mail,
d'une facture adressée à un client, d'un contrat à signer ou d'un message
enrichi sur mobile, chaque point de contact véhicule l'image de marque de
l'entreprise et participe à la relation qu'elle entretient avec ses
publics. La multiplication des canaux numériques a rendu cette
communication à la fois plus riche et plus exigeante : les entreprises
doivent produire, rapidement et à grande échelle, des supports cohérents,
personnalisés et fidèles à leur identité visuelle.

Or la fabrication de ces supports demeure, dans bien des cas, un exercice
lent et technique. Concevoir un e-mail responsive correct suppose de
maîtriser des langages de balisage tels que le HTML ou le MJML~; produire
une facture ou un contrat exige de la rigueur documentaire~; et l'ensemble
requiert le plus souvent l'intervention conjointe de profils créatifs et
techniques. Les équipes métier, en particulier les équipes marketing,
souhaitent au contraire pouvoir \emph{décrire} ce qu'elles veulent — ou
\emph{réutiliser} un visuel existant — et obtenir instantanément un
résultat professionnel et immédiatement modifiable.

C'est dans ce contexte que s'inscrit le présent projet de fin d'études,
réalisé au sein de \textbf{France Téléphone}, filiale du groupe
\textbf{Winvest Capital}. Il a consisté à concevoir et à développer
\textbf{Winaity Template Builder}, une plateforme logicielle en mode SaaS
(\emph{Software as a Service}) multi-entreprises destinée à la création, à
la gestion et à la génération de modèles de communication sur trois canaux :
l'e-mail, la facture et le contrat. Le c\oe{}ur de
la plateforme est un éditeur visuel par glisser-déposer, augmenté d'un
assistant d'intelligence artificielle \emph{agentique} capable de générer
et de modifier de véritables blocs éditables, voire de reconstruire une
campagne complète à partir d'une simple affiche marketing.

L'avènement récent des grands modèles de langage a ouvert des perspectives
inédites en matière d'assistance à la création de contenu. Le projet
explore précisément cette convergence entre l'ingénierie logicielle et
l'intelligence artificielle, en s'appuyant sur un modèle de langage
auto-hébergé afin de préserver la souveraineté des données de l'entreprise.

Ce rapport rend compte de la démarche suivie et des réalisations
accomplies. Il s'organise en six chapitres.

Le \textbf{premier chapitre} présente le cadre général du projet :
l'organisme d'accueil, la problématique, une étude critique des solutions
existantes, la solution proposée et la méthodologie de travail retenue.

Le \textbf{deuxième chapitre} est consacré à l'analyse des besoins et à la
conception : identification des acteurs, expression des besoins
fonctionnels et non fonctionnels, diagrammes de conception, architecture de
l'application, backlog produit et environnement de développement.

Les \textbf{quatre chapitres suivants} détaillent le déroulement des quatre
sprints du projet. Le troisième chapitre traite de l'authentification, de
l'isolation multi-entreprises et de la gestion des utilisateurs. Le
quatrième porte sur l'éditeur de modèles multi-canal. Le cinquième, c\oe{}ur
de la valeur ajoutée du projet, décrit l'assistant d'intelligence
artificielle agentique. Le sixième aborde les abonnements, les
intégrations et la mise en production.

Enfin, une \textbf{conclusion générale} dresse le bilan du travail réalisé,
des compétences acquises et des perspectives d'évolution de la plateforme.

% ---------------------------------------------------------------------
%  CHAPITRE 1
% ---------------------------------------------------------------------
\chapter{Contexte général du projet}
\label{chap:contexte}

\section*{Introduction}
Ce premier chapitre pose le cadre du projet. Nous présentons d'abord
l'organisme d'accueil au sein duquel le stage a été réalisé, puis le
projet lui-même : sa problématique, une étude critique des solutions
existantes sur le marché et la solution que nous proposons. Nous concluons
par la présentation et la justification de la méthodologie de travail
adoptée.

\section{Présentation de l'organisme d'accueil}

\subsection{Le groupe Winvest Capital}
Le présent projet de fin d'études a été réalisé au sein de
\textbf{France Téléphone}, entité rattachée au groupe
\textbf{Winvest Capital}. Winvest Capital est un groupe
[À COMPLÉTER — secteur / activité principale du groupe] regroupant plusieurs
entités opérationnelles, dont France Téléphone.

\figplaceholder{Logo du groupe Winvest Capital}{fig:logo-winvest}{3cm}

\begin{description}[leftmargin=2.2cm,style=nextline,font=\bfseries]
  \item[Raison sociale] [À COMPLÉTER]
  \item[Date de création] [À COMPLÉTER]
  \item[Siège social] [À COMPLÉTER]
  \item[Secteur d'activité] [À COMPLÉTER]
  \item[Effectif] [À COMPLÉTER]
\end{description}

\subsection{France Téléphone}
\textbf{France Téléphone} est l'entité du groupe au sein de laquelle s'est
déroulé le stage. Son activité s'inscrit dans le domaine
[À COMPLÉTER — ex. de la distribution de forfaits et de services de
télécommunication], ce qui l'amène à produire de manière régulière des
supports de communication commerciale : offres promotionnelles, campagnes
e-mail, factures, contrats d'abonnement et messages à destination de sa
clientèle.

\figplaceholder{Logo de France Téléphone}{fig:logo-ft}{3cm}

Cette intensité de communication commerciale, sur des canaux multiples et
à un rythme soutenu, constitue précisément le terrain qui a fait émerger le
besoin à l'origine de ce projet.

\subsection{Cadre du stage}
Le stage s'est déroulé sur une durée de six mois, du 2 février au
31 juillet 2026, au sein de l'équipe [À COMPLÉTER — ex. développement /
systèmes d'information]. La mission confiée consistait à concevoir et
développer,
de bout en bout, une plateforme interne et commercialisable de création de
modèles de communication assistée par intelligence artificielle. Le travail
a couvert l'ensemble du cycle de vie logiciel : analyse des besoins,
conception, développement, tests et mise en production.

\section{Présentation du projet}

\subsection{Problématique}
La production de supports de communication professionnels se heurte à
plusieurs difficultés récurrentes que l'on peut synthétiser ainsi.

\begin{itemize}[leftmargin=1.4cm]
  \item \textbf{Une fabrication lente et technique.} La création d'un
  e-mail responsive de qualité suppose la maîtrise du HTML ou du MJML et
  une bonne connaissance des contraintes de compatibilité entre clients de
  messagerie. Les documents structurés — factures, contrats — imposent en
  outre une mise en forme rigoureuse et paginée.

  \item \textbf{Une dépendance à des compétences croisées.} Obtenir un
  résultat à la fois esthétique et techniquement correct requiert souvent
  la collaboration de profils créatifs (design) et de profils techniques
  (intégration), ce qui allonge les délais et alourdit le processus.

  \item \textbf{La multiplicité des canaux.} Une même campagne peut devoir
  associer un e-mail, une facture et un contrat. Or les outils du marché
  sont le plus souvent spécialisés sur un seul canal.

  \item \textbf{Le cloisonnement des données.} Dans un contexte
  multi-entreprises, chaque société doit disposer d'un espace de travail
  strictement isolé : ses modèles, ses utilisateurs et son identité de
  marque ne doivent jamais être visibles par une autre.

  \item \textbf{La rigidité des assistants existants.} Les générateurs
  automatiques disponibles produisent généralement du code HTML figé,
  difficile à reprendre et à modifier finement, plutôt que des éléments
  structurés directement éditables.
\end{itemize}

La problématique centrale du projet peut dès lors se formuler ainsi :
\emph{comment permettre à une équipe métier, sans compétence technique
particulière, de produire rapidement des supports de communication
multi-canaux, cohérents avec l'identité de son entreprise, tout en
garantissant l'isolation des données propre à un contexte
multi-entreprises et en tirant parti de l'intelligence artificielle sans
compromettre la souveraineté de ces données~?}

\subsection{Étude de l'existant}
Afin de situer notre solution, nous avons analysé plusieurs catégories
d'outils existants représentatifs du marché.

\begin{itemize}[leftmargin=1.4cm]
  \item \textbf{Canva} — Outil de conception graphique généraliste très
  répandu. Il excelle pour produire des visuels, mais n'est pas pensé pour
  l'e-mail responsive (pas de sortie MJML/HTML native robuste), ne couvre
  pas les documents transactionnels (facture, contrat) ni la logique
  multi-entreprises, et son assistance IA reste orientée image plutôt que
  génération de structures éditables par canal.

  \item \textbf{Mailchimp} — Plateforme d'e-mailing reconnue, dotée d'un
  éditeur par blocs et de fonctions marketing. Elle demeure toutefois
  centrée sur un seul canal (l'e-mail), pensée pour un compte unique plutôt
  que pour une architecture multi-locataire déléguée, et ne propose pas de
  reconstruction de campagne à partir d'une image.

  \item \textbf{Stripo et Beefree} — Éditeurs d'e-mails par glisser-déposer
  produisant du HTML/MJML de bonne qualité et exportables. Ils offrent une
  expérience visuelle soignée, mais ne fournissent pas d'assistant
  \emph{agentique} générant de vrais blocs éditables, ne couvrent pas
  l'ensemble des canaux (facture, contrat) et n'intègrent pas de
  fonction «~affiche vers e-mail~».

  \item \textbf{Éditeurs MJML bruts} — Des outils tels que l'éditeur
  officiel MJML s'adressent à des développeurs et supposent l'écriture
  directe de balises~; ils sont donc inaccessibles à un utilisateur métier.

  \item \textbf{Générateurs d'e-mails par IA} — Certains services génèrent
  un e-mail à partir d'une invite en langage naturel, mais restituent un
  bloc de HTML monolithique, difficile à retoucher, sans modèle de données
  structuré ni contrôle fin du rendu.
\end{itemize}

\subsection{Critique de l'existant}
De cette étude se dégage un constat clair : aucun des outils analysés ne
réunit simultanément l'ensemble des propriétés requises par notre contexte.
Le tableau~\ref{tab:comparatif-existant} synthétise cette comparaison.

\begin{table}[H]
\centering
\caption{Comparaison synthétique des solutions existantes}
\label{tab:comparatif-existant}
\footnotesize
\setlength{\tabcolsep}{4pt}
\begin{tabularx}{\linewidth}{Xccccc}
\toprule
\textbf{Critère} & \textbf{Canva} & \textbf{Mailchimp} & \textbf{Stripo /}
& \textbf{Éditeurs} & \textbf{Winaity} \\
 & & & \textbf{Beefree} & \textbf{MJML} & \\
\midrule
Éditeur visuel par blocs        & \checkmark & \checkmark & \checkmark & --         & \checkmark \\
Multi-canal (3 types)           & --         & --         & --         & --         & \checkmark \\
Multi-entreprises (tenants)     & --         & partiel    & --         & --         & \checkmark \\
IA agentique (blocs éditables)  & --         & partiel    & --         & --         & \checkmark \\
Affiche $\rightarrow$ e-mail    & --         & --         & --         & --         & \checkmark \\
LLM auto-hébergé (souveraineté) & --         & --         & --         & --         & \checkmark \\
Accessible sans compétence tech.& \checkmark & \checkmark & \checkmark & --         & \checkmark \\
\bottomrule
\end{tabularx}
\end{table}

Les principales lacunes identifiées sont les suivantes~:
\begin{itemize}[leftmargin=1.4cm]
  \item une \textbf{spécialisation mono-canal} qui oblige à multiplier les
  outils~;
  \item l'\textbf{absence d'un vrai modèle de blocs éditables} partagé
  entre l'édition manuelle et l'édition assistée par IA~;
  \item l'\textbf{absence d'isolation multi-entreprises déléguée} adaptée à
  un éditeur de logiciels souhaitant l'intégrer à ses propres outils~;
  \item la \textbf{dépendance à des API d'IA tierces}, incompatible avec les
  exigences de confidentialité des données d'une entreprise.
\end{itemize}

\subsection{Solution proposée}
Pour répondre à ces limites, nous proposons \textbf{Winaity Template
Builder} : une plateforme SaaS multi-entreprises de création de modèles de
communication, articulée autour de trois idées directrices.

\begin{enumerate}[leftmargin=1.4cm]
  \item \textbf{Un éditeur visuel multi-canal unique.} Trois types de
  modèles — e-mail, facture et contrat — sont créés depuis une
  même plateforme, chacun disposant d'un éditeur dédié. L'éditeur d'e-mail,
  produit phare, repose sur un canevas de blocs manipulés par
  glisser-déposer et sur un modèle de données structuré.

  \item \textbf{Un assistant d'IA agentique produisant de vrais blocs.} Au
  lieu de générer du HTML jetable, l'assistant appelle des \emph{outils} de
  construction qui renvoient exactement le même format de blocs que
  l'éditeur manuel. Édition visuelle et édition par IA sont donc
  parfaitement interchangeables. L'assistant sait générer un e-mail complet
  à partir d'une consigne, le modifier en conversation, ou reconstruire une
  campagne à partir d'une \emph{affiche} lue par vision par ordinateur.

  \item \textbf{Une plateforme d'entreprise souveraine.} L'ensemble des
  données est cloisonné par entreprise (\emph{tenant}). Le modèle de langage
  qui alimente l'IA est \textbf{auto-hébergé}, ce qui garantit que les
  données ne quittent jamais l'infrastructure du groupe. La plateforme
  intègre par ailleurs la gestion des rôles, un système d'abonnements et
  une couche d'intégration permettant à des outils externes de l'embarquer.
\end{enumerate}

\section{Méthodologie de travail}

\subsection{Choix d'une méthodologie agile}
Le projet se caractérise par un périmètre fonctionnel large, des exigences
susceptibles d'évoluer au fil des retours de l'entreprise, et une forte
composante d'exploration technique (notamment sur le volet intelligence
artificielle). Ces caractéristiques rendent inadaptées les approches de
développement séquentielles, dites «~en cascade~», qui figent les besoins
en début de projet et n'autorisent que tardivement la validation par le
client. Nous avons donc opté pour une \textbf{approche agile}, qui privilégie
les livraisons incrémentales, les retours fréquents et l'adaptation
continue. Le tableau~\ref{tab:agile-vs-cascade} résume cette comparaison.

\begin{table}[H]
\centering
\caption{Approche agile versus approche traditionnelle}
\label{tab:agile-vs-cascade}
\small
\begin{tabularx}{\linewidth}{lXX}
\toprule
\textbf{Aspect} & \textbf{Approche agile} & \textbf{Approche en cascade} \\
\midrule
Cycle de développement & Itératif et incrémental & Séquentiel et linéaire \\
Gestion du changement & Adaptation continue & Faible tolérance au changement \\
Livraison & Incréments fréquents & Livraison en fin de projet \\
Implication du client & Retours réguliers & Implication tardive \\
Gestion du risque & Réduite par itérations & Concentrée en fin de cycle \\
Documentation & Progressive & Exhaustive en amont \\
\bottomrule
\end{tabularx}
\end{table}

\subsection{Choix de Scrum}
Parmi les méthodes agiles (Scrum, Extreme Programming, Kanban), nous avons
retenu \textbf{Scrum} pour la clarté de son cadre organisationnel et son
adéquation à un projet mené sur plusieurs mois. Scrum structure le travail
en itérations de durée fixe appelées \emph{sprints}, à l'issue desquelles
un incrément potentiellement livrable est produit. Il définit trois rôles :

\begin{itemize}[leftmargin=1.4cm]
  \item le \textbf{Product Owner}, qui porte la vision du produit et
  priorise le \emph{backlog}~;
  \item le \textbf{Scrum Master}, qui veille au bon déroulement du cadre et
  lève les obstacles~;
  \item l'\textbf{équipe de développement}, qui réalise l'incrément.
\end{itemize}

Dans le contexte particulier de ce stage, ces rôles ont été adaptés à une
configuration réduite : l'encadrant professionnel a joué le rôle de Product
Owner en exprimant et en priorisant les besoins, tandis que la réalisation
technique a été assurée dans le cadre du stage. Des points d'avancement
réguliers ont tenu lieu de cérémonies Scrum (planification, revue de
sprint, rétrospective).

\figplaceholder{Le processus Scrum}{fig:scrum}{5cm}

Le projet a été découpé en \textbf{quatre sprints}, chacun correspondant à
un ensemble cohérent de fonctionnalités~: (1) authentification, isolation
multi-entreprises et gestion des utilisateurs~; (2) éditeur de modèles
multi-canal~; (3) assistant d'intelligence artificielle agentique~; (4)
abonnements, intégrations et mise en production. Ce découpage est détaillé
au chapitre suivant.

\section*{Conclusion}
Ce chapitre a introduit l'organisme d'accueil, France Téléphone au sein du
groupe Winvest Capital, puis a posé la problématique du projet à travers
une étude critique des solutions existantes, faisant ressortir l'absence
d'un outil combinant approche multi-canal, isolation multi-entreprises,
intelligence artificielle agentique et souveraineté des données. Nous avons
enfin justifié le choix de la méthodologie Scrum. Le chapitre suivant
détaille l'analyse des besoins et la conception de la plateforme.

% ---------------------------------------------------------------------
%  CHAPITRE 2
% ---------------------------------------------------------------------
\chapter{Analyse des besoins et conception}
\label{chap:analyse}

\section*{Introduction}
Après avoir situé le projet, ce chapitre en établit les fondations. Nous
identifions d'abord les acteurs de la plateforme et exprimons les besoins
fonctionnels et non fonctionnels. Nous présentons ensuite la conception à
travers le diagramme de cas d'utilisation global, le diagramme de classes
et l'architecture applicative. Nous terminons par l'initialisation du
projet : le backlog produit, la planification des sprints et
l'environnement de développement.

\section{Analyse des besoins}

\subsection{Identification des acteurs}
Un acteur désigne un rôle joué par une entité externe qui interagit avec le
système. Pour Winaity Template Builder, nous distinguons les acteurs
suivants.

\begin{itemize}[leftmargin=1.4cm]
  \item \textbf{Utilisateur / Éditeur.} Membre d'une entreprise cliente. Il
  crée, modifie, prévisualise et exporte des modèles, utilise l'assistant
  d'IA et organise ses modèles (favoris, duplication).

  \item \textbf{Administrateur d'entreprise (Admin).} Utilisateur disposant
  de droits étendus au sein de son entreprise. Outre les capacités d'un
  éditeur, il invite et gère les membres de son équipe, souscrit un
  abonnement et gère la facturation, et génère des clés d'intégration
  (API).

  \item \textbf{Membre marketing.} Rôle spécialisé se comportant comme un
  éditeur, mais habilité en plus à constituer et maintenir la galerie de
  modèles prédéfinis propre à son entreprise, par catégorie.

  \item \textbf{Super-administrateur.} Acteur interne à l'éditeur de la
  plateforme. Il supervise l'ensemble des entreprises, consulte les
  abonnements et attribue les plans, y compris le plan interne réservé aux
  entités du groupe.

  \item \textbf{Client applicatif (M2M).} Système externe (par exemple un
  CRM) authentifié par une clé d'API, qui embarque l'éditeur pour ses
  propres utilisateurs via le flux d'intégration.
\end{itemize}

\subsection{Besoins fonctionnels}
Les besoins fonctionnels décrivent les services que la plateforme doit
rendre. Ils sont regroupés ci-dessous par domaine.

\paragraph{Authentification et gestion des comptes.}
\begin{itemize}[leftmargin=1.4cm,nosep]
  \item S'inscrire et créer une entreprise, ou rejoindre une entreprise
  existante sur invitation.
  \item S'authentifier de manière sécurisée (jeton JWT) et gérer sa
  session.
  \item Consulter et modifier son profil.
\end{itemize}

\paragraph{Multi-entreprises et gestion des utilisateurs.}
\begin{itemize}[leftmargin=1.4cm,nosep]
  \item Garantir l'isolation stricte des données par entreprise.
  \item Inviter des membres, accepter une invitation, lister l'équipe.
  \item Attribuer des rôles (administrateur, éditeur, marketing).
\end{itemize}

\paragraph{Édition de modèles multi-canal.}
\begin{itemize}[leftmargin=1.4cm,nosep]
  \item Créer, modifier, dupliquer, supprimer et lister des modèles.
  \item Éditer un e-mail par glisser-déposer de blocs, prévisualiser en
  direct et éditer le code.
  \item Créer des factures et des contrats.
  \item Exporter en HTML responsive et en PDF, substituer des variables de
  personnalisation.
  \item Gérer les favoris et la galerie de modèles prédéfinis.
\end{itemize}

\paragraph{Assistance par intelligence artificielle.}
\begin{itemize}[leftmargin=1.4cm,nosep]
  \item Générer un e-mail complet à partir d'une consigne en langage
  naturel.
  \item Modifier un modèle en conversation (thème, image, disposition,
  ajout ou déplacement de blocs).
  \item Appliquer une modification à un élément précis sélectionné dans le
  canevas.
  \item Reconstruire un e-mail éditable à partir d'une affiche marketing.
\end{itemize}

\paragraph{Abonnements, intégrations et administration.}
\begin{itemize}[leftmargin=1.4cm,nosep]
  \item Souscrire, changer, annuler et réactiver un abonnement~; appliquer
  les quotas par plan.
  \item Générer des clés d'API et embarquer l'éditeur depuis un outil
  externe.
  \item Superviser les entreprises et attribuer les plans (super-admin).
\end{itemize}

\subsection{Besoins non fonctionnels}
Les besoins non fonctionnels définissent les contraintes de qualité que la
plateforme doit respecter.

\begin{itemize}[leftmargin=1.4cm]
  \item \textbf{Sécurité et confidentialité.} Authentification par jeton,
  cloisonnement rigoureux des données par \texttt{tenant\_id}, contrôle
  d'accès basé sur les rôles, secrets tenus hors du dépôt de code, et
  \textbf{souveraineté des données} garantie par un modèle de langage
  auto-hébergé.
  \item \textbf{Performance.} Génération d'un e-mail par l'IA en une dizaine
  de secondes, rendu et prévisualisation réactifs.
  \item \textbf{Évolutivité (scalabilité).} Architecture microservices
  permettant de faire évoluer et de déployer chaque service
  indépendamment.
  \item \textbf{Maintenabilité.} Séparation des responsabilités, patron
  CQRS côté services, modèle de données partagé entre édition manuelle et
  édition IA.
  \item \textbf{Utilisabilité.} Interface intuitive, accessible à un
  utilisateur métier sans compétence technique.
  \item \textbf{Fiabilité.} Confirmations honnêtes de l'assistant (aucun
  succès annoncé si le modèle n'a pas réellement changé), garde de
  contraste garantissant la lisibilité, gestion propre des erreurs.
\end{itemize}

\section{Conception}

\subsection{Diagramme de cas d'utilisation global}
Le diagramme de cas d'utilisation offre une vue d'ensemble du comportement
fonctionnel de la plateforme et des interactions entre acteurs et cas
d'utilisation. La figure~\ref{fig:uc-global} présente le diagramme global du
système.

\figplaceholder{Diagramme de cas d'utilisation global}{fig:uc-global}{8cm}

\subsection{Diagramme de classes}
Le diagramme de classes décrit la structure statique du système : les
entités, leurs attributs et leurs relations. Les principales entités du
domaine sont \texttt{Tenant} (entreprise), \texttt{User} (utilisateur et
son rôle), \texttt{Template} (modèle, caractérisé par son type de canal et
son contenu), les entités d'invitation et de jeton, ainsi que les entités
liées à l'abonnement (plan, cycle de facturation, statut). La
figure~\ref{fig:classes} présente le diagramme de classes.

\figplaceholder{Diagramme de classes}{fig:classes}{9cm}

\subsection{Architecture de l'application}
La plateforme adopte une \textbf{architecture microservices}. Cette section
en présente les vues physique et logique.

\subsubsection{Architecture physique}
L'architecture physique décrit les composants déployés et leurs
interactions. Elle se compose des éléments suivants~:

\begin{itemize}[leftmargin=1.4cm]
  \item \textbf{Client (navigateur).} Interface utilisateur exécutée dans le
  navigateur.
  \item \textbf{Application frontend.} Application Next.js servant
  l'interface et hébergeant les routes d'IA côté serveur.
  \item \textbf{Passerelle API (API Gateway).} Unique surface REST publique~;
  elle traduit les appels REST en appels gRPC vers les services internes et
  gère l'authentification, la facturation et l'application des quotas.
  \item \textbf{Service d'authentification.} Gère les comptes, les
  entreprises, les rôles et les abonnements.
  \item \textbf{Service de modèles.} Gère le cycle de vie des modèles
  (CRUD, rendu, duplication, favoris).
  \item \textbf{Services d'IA et d'images.} Services auxiliaires
  (suggestions de palettes, mapping de variables, recherche d'images) et
  \textbf{serveur de modèle de langage auto-hébergé}.
  \item \textbf{Bases de données PostgreSQL} (\texttt{auth\_db} et
  \texttt{templates\_db}) et \textbf{stockage objet MinIO} pour les médias.
\end{itemize}

\figplaceholder{Architecture physique de l'application}{fig:archi-phys}{7cm}

\subsubsection{Architecture logique}
Sur le plan logique, la plateforme suit un modèle en couches communiquant
selon un schéma requête/réponse. Le client émet une requête REST
(authentifiée par un jeton JWT porté par un \emph{cookie}) vers la
passerelle API. Celle-ci relaie la requête, via \textbf{gRPC}, au service
concerné. Les services, développés avec le framework NestJS, appliquent le
patron \textbf{CQRS} (séparation des commandes et des requêtes) et
persistent leurs données au moyen de l'ORM TypeORM. La
figure~\ref{fig:archi-log} illustre cette architecture.

\figplaceholder{Architecture logique de l'application}{fig:archi-log}{7cm}

Le choix du \textbf{gRPC} pour la communication inter-services se justifie
par ses performances (sérialisation binaire via Protocol Buffers) et par le
contrat d'interface fort qu'il impose. Le jeton JWT transporte la charge
utile \texttt{\{ userId, tenantId, email, role \}}, où le
\texttt{tenant\_id} constitue la clé d'isolation présente dans tous les
services.

\section{Initialisation du projet}

\subsection{Backlog produit}
Le backlog produit recense les fonctionnalités attendues, exprimées sous
forme d'\emph{epics} et priorisées. Le tableau~\ref{tab:backlog} en présente
la synthèse.

\begin{table}[H]
\centering
\caption{Backlog produit (synthèse par epic)}
\label{tab:backlog}
\small
\begin{tabularx}{\linewidth}{clXc}
\toprule
\textbf{ID} & \textbf{Epic} & \textbf{Description} & \textbf{Priorité} \\
\midrule
1  & Authentification & Inscription, connexion sécurisée (JWT), profil. & Haute \\
2  & Multi-entreprises & Isolation des données par \texttt{tenant\_id}. & Haute \\
3  & Gestion des utilisateurs & Invitations, équipe, rôles (admin / éditeur / marketing), super-admin. & Haute \\
4  & Éditeur d'e-mail & Canevas de blocs, glisser-déposer, aperçu, code, export HTML/PDF. & Haute \\
5  & Modèles multi-canal & Facture et contrat. & Haute \\
6  & Bibliothèque de modèles & Favoris, duplication, galerie de modèles prédéfinis. & Moyenne \\
7  & Génération IA par consigne & E-mail complet à partir d'une invite en langage naturel. & Haute \\
8  & Édition IA conversationnelle & Modification par dialogue et édition ciblée par sélection. & Haute \\
9  & Affiche $\rightarrow$ e-mail & Reconstruction d'un e-mail éditable à partir d'une image. & Haute \\
10 & MCP \& LLM auto-hébergé & Exposition des outils via MCP, inférence sur serveur vLLM privé. & Moyenne \\
11 & Abonnements et facturation & Plans, paiement Stripe, TVA, quotas par plan. & Haute \\
12 & Super-administration & Supervision des entreprises, attribution des plans. & Moyenne \\
13 & Intégrations et embed & Clés d'API, embarquement de l'éditeur par un outil externe. & Moyenne \\
14 & Site marketing & Site public de présentation (accueil, offres, contact\dots). & Basse \\
15 & Mise en production & Conteneurisation, déploiement, exposition sécurisée. & Haute \\
\bottomrule
\end{tabularx}
\end{table}

\subsection{Planification des sprints}
Le projet a été organisé en quatre sprints regroupant les epics du backlog
de manière cohérente, comme le résume le tableau~\ref{tab:sprints}.

\begin{table}[H]
\centering
\caption{Planification des sprints}
\label{tab:sprints}
\small
\begin{tabularx}{\linewidth}{lX}
\toprule
\textbf{Sprint} & \textbf{Contenu} \\
\midrule
Sprint 1 & Authentification, multi-entreprises et gestion des utilisateurs (epics 1--3). \\
Sprint 2 & Éditeur de modèles multi-canal et bibliothèque de modèles (epics 4--6). \\
Sprint 3 & Assistant d'intelligence artificielle agentique (epics 7--10). \\
Sprint 4 & Abonnements, intégrations et mise en production (epics 11--15). \\
\bottomrule
\end{tabularx}
\end{table}

La figure~\ref{fig:gantt} présente le diagramme de Gantt du projet, qui
positionne dans le temps la phase d'analyse et de conception, les quatre
sprints de réalisation et la rédaction du rapport, sur la durée du stage
(du 2 février au 31 juillet 2026).

\begin{figure}[H]
\centering
\resizebox{\linewidth}{!}{%
\begin{ganttchart}[
  hgrid, vgrid,
  x unit=2.2cm,
  y unit title=0.7cm,
  y unit chart=0.8cm,
  title label font=\bfseries\small,
  bar/.append style={fill=blue!30},
  bar height=.6,
  bar label font=\small
]{1}{6}
\gantttitle{Février}{1}\gantttitle{Mars}{1}\gantttitle{Avril}{1}\gantttitle{Mai}{1}\gantttitle{Juin}{1}\gantttitle{Juillet}{1}\\
\ganttbar{Analyse \& conception}{1}{1}\\
\ganttbar{Sprint 1 : Authentification \& multi-entreprises}{1}{2}\\
\ganttbar{Sprint 2 : Éditeur multi-canal}{2}{3}\\
\ganttbar{Sprint 3 : Assistant IA agentique}{4}{5}\\
\ganttbar{Sprint 4 : Abonnements \& mise en production}{5}{6}\\
\ganttbar{Rédaction du rapport}{6}{6}
\end{ganttchart}}
\caption{Diagramme de Gantt du projet (février--juillet 2026)}
\label{fig:gantt}
\end{figure}

\subsection{Environnement de développement}

\subsubsection{Environnement matériel}
Le développement s'est appuyé sur les ressources suivantes~:
\begin{itemize}[leftmargin=1.4cm,nosep]
  \item un poste de travail personnel~: [À COMPLÉTER — processeur, mémoire
  RAM, système d'exploitation]~;
  \item un \textbf{serveur d'entreprise doté d'un accélérateur graphique
  (GPU)} hébergeant le moteur d'inférence vLLM et servant également au
  déploiement en production.
\end{itemize}

\subsubsection{Environnement logiciel et technologies}
Le tableau~\ref{tab:techno} récapitule la pile technologique. Les
principaux choix sont ensuite justifiés.

\begin{table}[H]
\centering
\caption{Pile technologique du projet}
\label{tab:techno}
\small
\begin{tabularx}{\linewidth}{lX}
\toprule
\textbf{Couche} & \textbf{Technologies} \\
\midrule
Frontend & Next.js~16 (App Router), React~19, TypeScript, Tailwind CSS, Framer Motion \\
Édition & dnd-kit (glisser-déposer), Monaco (éditeur de code), MJML (rendu e-mail), Tiptap (texte riche) \\
Backend & NestJS, gRPC / Protocol Buffers, patron CQRS, TypeORM \\
Données & PostgreSQL (\texttt{auth\_db}, \texttt{templates\_db}), MinIO (stockage objet) \\
IA & Vercel AI SDK, MCP SDK, serveur vLLM auto-hébergé (Gemma / Qwen), services Python auxiliaires \\
Paiement & Stripe (Checkout intégré, abonnements, TVA) \\
Industrialisation & Docker \& Docker Compose, GitLab (CI), Cloudflare Tunnel \\
Outils & Visual Studio Code, Git, Postman, PlantUML / draw.io \\
\bottomrule
\end{tabularx}
\end{table}

\paragraph{Frontend — Next.js et React.}
Next.js, fondé sur React, a été retenu pour son \emph{App Router}, sa
capacité à exécuter du code côté serveur (indispensable pour héberger en
toute sécurité les appels au modèle de langage) et son rendu performant.
React apporte son modèle de composants et son écosystème mature~[1,~2].

\figplaceholder{Next.js et React}{fig:logo-next-react}{2.5cm}

\paragraph{Backend — NestJS et gRPC.}
NestJS structure le code serveur autour de modules, contrôleurs et
services, et facilite l'application du patron CQRS. La communication
inter-services repose sur gRPC, qui offre un contrat d'interface fort et une
sérialisation binaire efficace via Protocol Buffers~[3,~4].

\paragraph{Données — PostgreSQL, TypeORM et MinIO.}
PostgreSQL assure la persistance relationnelle des deux bases du système.
TypeORM fait le pont entre les entités du domaine et la base. MinIO,
compatible avec l'API S3, stocke les médias (images) de la
plateforme~[5,~6,~7].

\paragraph{Rendu e-mail — MJML.}
MJML est un langage de balisage qui compile vers du HTML d'e-mail
responsive compatible avec la majorité des clients de messagerie~; il
constitue le format d'export de l'éditeur d'e-mail~[8].

\paragraph{Intelligence artificielle — Vercel AI SDK, MCP et vLLM.}
Le Vercel AI SDK orchestre la boucle agentique d'appel d'outils côté
serveur. Le Model Context Protocol (MCP) expose les outils de construction
au modèle de manière standardisée. L'inférence est assurée par un serveur
\textbf{vLLM auto-hébergé} exécutant le modèle Gemma (et Qwen en repli),
tous deux compatibles avec l'API OpenAI, ce qui garantit que les données ne
quittent pas l'infrastructure du groupe~[9,~10,~11,~12].

\paragraph{Paiement — Stripe.}
Stripe gère les abonnements récurrents, le Checkout intégré et le calcul de
la TVA~[13].

\paragraph{Industrialisation — Docker, GitLab et Cloudflare.}
Docker et Docker Compose conteneurisent l'ensemble des services et
garantissent la reproductibilité des environnements. Le dépôt est hébergé
sur GitLab. L'exposition publique en production s'appuie sur un tunnel
Cloudflare, sécurisé et sans ouverture de port entrant~[14,~15].

\section*{Conclusion}
Ce chapitre a précisé les acteurs et les besoins de la plateforme, présenté
sa conception (cas d'utilisation, classes, architecture microservices) et
posé les bases de sa réalisation : backlog produit, planification en quatre
sprints et environnement technique. Les chapitres suivants détaillent, sprint
par sprint, la mise en \oe{}uvre de ces choix.

% ---------------------------------------------------------------------
%  CHAPITRE 3
% ---------------------------------------------------------------------
\chapter{Sprint~1 : Authentification, multi-entreprises et gestion des utilisateurs}
\label{chap:sprint1}

\section*{Introduction}
Ce premier sprint pose les fondations de la plateforme : la manière dont un
utilisateur accède au système, dont ses données sont cloisonnées par
entreprise, et dont les membres d'une même organisation sont gérés. Nous
présentons successivement les objectifs du sprint, le raffinement des cas
d'utilisation concernés, les diagrammes de séquence associés, puis la
réalisation obtenue.

\section{Objectifs du sprint}
Ce sprint couvre les epics~1 à~3 du backlog produit. Le
tableau~\ref{tab:backlog-s1} en détaille les user stories.

\begin{table}[H]
\centering
\caption{Backlog du sprint~1}
\label{tab:backlog-s1}
\small
\begin{tabularx}{\linewidth}{clX}
\toprule
\textbf{ID} & \textbf{User story} & \textbf{Description} \\
\midrule
1.1 & S'inscrire & En tant qu'utilisateur, je veux créer un compte et une entreprise afin de disposer d'un espace de travail. \\
1.2 & Se connecter & En tant qu'utilisateur, je veux m'authentifier de façon sécurisée pour accéder à la plateforme. \\
1.3 & Gérer mon profil & En tant qu'utilisateur, je veux consulter et modifier mes informations personnelles. \\
2.1 & Isolation des données & En tant qu'entreprise, je veux que mes modèles et mes membres soient invisibles pour les autres entreprises. \\
3.1 & Inviter un membre & En tant qu'administrateur, je veux inviter des collaborateurs à rejoindre mon entreprise. \\
3.2 & Rejoindre une entreprise & En tant qu'invité, je veux accepter une invitation pour intégrer une entreprise. \\
3.3 & Gérer l'équipe & En tant qu'administrateur, je veux lister les membres et leur attribuer un rôle (administrateur, éditeur, marketing). \\
3.4 & Superviser les entreprises & En tant que super-administrateur, je veux consulter l'ensemble des entreprises et de leurs abonnements. \\
\bottomrule
\end{tabularx}
\end{table}

\section{Raffinement et description des cas d'utilisation}
Nous précisons dans cette section le fonctionnement des services
implémentés, en détaillant les processus clés et les interactions entre
acteurs et système.

\subsection{Cas d'utilisation « Authentification et inscription »}
La figure~\ref{fig:uc-auth} illustre les cas d'utilisation relatifs à
l'inscription et à l'authentification. L'inscription crée conjointement un
compte utilisateur et l'entreprise associée~; la connexion s'effectue par
identifiants et délivre un jeton de session.

\figplaceholder{Cas d'utilisation « Authentification et inscription »}{fig:uc-auth}{7cm}

Le tableau~\ref{tab:uc-inscription} décrit textuellement le scénario
« S'inscrire ».

\begin{table}[H]
\centering
\caption{Description textuelle du cas d'utilisation « S'inscrire »}
\label{tab:uc-inscription}
\small
\begin{tabularx}{\linewidth}{lX}
\toprule
\textbf{Cas d'utilisation} & S'inscrire \\
\textbf{Acteur} & Utilisateur \\
\textbf{Préconditions} & L'utilisateur n'est pas connecté. \\
\textbf{Postconditions} & Un compte et une entreprise sont créés~; l'utilisateur est authentifié. \\
\midrule
\textbf{Scénario nominal} &
1.~L'utilisateur accède au formulaire d'inscription.\newline
2.~Il saisit ses informations et le nom de son entreprise.\newline
3.~Le système valide les données fournies.\newline
4.~Le système crée l'entreprise puis le compte associé.\newline
5.~Le système génère un jeton de session et redirige l'utilisateur vers le tableau de bord. \\
\midrule
\textbf{Scénario alternatif} &
3.a.~Si l'adresse e-mail ou le nom d'entreprise existe déjà, le système en informe l'utilisateur et l'invite à corriger sa saisie. \\
\bottomrule
\end{tabularx}
\end{table}

\subsection{Cas d'utilisation « Gestion des utilisateurs »}
La figure~\ref{fig:uc-users} présente les cas d'utilisation relatifs à la
gestion de l'équipe : invitation d'un membre, acceptation d'une invitation,
consultation des membres et attribution des rôles.

\figplaceholder{Cas d'utilisation « Gestion des utilisateurs »}{fig:uc-users}{7cm}

Le tableau~\ref{tab:uc-invite} décrit le scénario « Inviter un membre ».

\begin{table}[H]
\centering
\caption{Description textuelle du cas d'utilisation « Inviter un membre »}
\label{tab:uc-invite}
\small
\begin{tabularx}{\linewidth}{lX}
\toprule
\textbf{Cas d'utilisation} & Inviter un membre \\
\textbf{Acteur} & Administrateur d'entreprise \\
\textbf{Préconditions} & L'administrateur est connecté. \\
\textbf{Postconditions} & Une invitation est émise et le futur membre est notifié. \\
\midrule
\textbf{Scénario nominal} &
1.~L'administrateur ouvre la page de gestion de l'équipe.\newline
2.~Il saisit l'adresse e-mail de l'invité et sélectionne un rôle.\newline
3.~Le système crée une invitation rattachée à l'entreprise et envoie un lien à l'invité.\newline
4.~L'invité suit le lien, définit son mot de passe et rejoint l'entreprise. \\
\midrule
\textbf{Scénario alternatif} &
2.a.~Si l'adresse correspond déjà à un membre de l'entreprise, le système signale que l'utilisateur en fait partie. \\
\bottomrule
\end{tabularx}
\end{table}

\section{Diagrammes de séquence}

\subsection{Séquence « S'inscrire »}
La figure~\ref{fig:seq-inscription} décrit les échanges entre le client, la
passerelle API et le service d'authentification lors d'une inscription. La
passerelle traduit la requête REST en appel gRPC~; le service valide les
données, crée l'entreprise et le compte, puis renvoie un jeton porté par un
\emph{cookie} sécurisé.

\figplaceholder{Diagramme de séquence « S'inscrire »}{fig:seq-inscription}{7cm}

\subsection{Séquence « Inviter un membre »}
La figure~\ref{fig:seq-invite} illustre le flux d'invitation d'un membre,
depuis la demande de l'administrateur jusqu'à l'acceptation par l'invité.

\figplaceholder{Diagramme de séquence « Inviter un membre »}{fig:seq-invite}{7cm}

\section{Réalisation}
Cette section présente les résultats obtenus au cours du sprint, tant sur le
plan technique que du point de vue des interfaces.

\subsection{Service d'authentification et modèle multi-entreprises}
Le service d'authentification a été développé avec NestJS selon le patron
CQRS et expose ses opérations en gRPC. Il prend en charge l'inscription, la
connexion, la validation de jeton, l'invitation et l'acceptation de membres,
la consultation et la mise à jour du profil, ainsi que le listage des
membres d'une entreprise. Les données sont persistées dans la base
\texttt{auth\_db} via l'ORM TypeORM.

Le c\oe{}ur du modèle multi-entreprises repose sur la notion de
\emph{tenant} (entreprise). Chaque compte est rattaché à un
\texttt{tenant\_id}, qui est embarqué dans la charge utile du jeton JWT aux
côtés de l'identifiant utilisateur, de l'adresse e-mail et du rôle. Ce
\texttt{tenant\_id} est systématiquement propagé aux services applicatifs
et sert de filtre d'isolation : aucune requête ne peut accéder aux données
d'une autre entreprise. La passerelle API constitue l'unique surface REST
publique et vérifie le jeton, transporté par un \emph{cookie}, avant de
relayer chaque appel en gRPC vers le service concerné.

\subsection{Inscription et connexion}
Les figures~\ref{fig:ui-register} et~\ref{fig:ui-login} présentent
respectivement les interfaces d'inscription et de connexion. L'inscription
recueille les informations de l'utilisateur ainsi que le nom de son
entreprise, puis crée l'espace de travail correspondant. La connexion
authentifie l'utilisateur et ouvre une session sécurisée.

\figplaceholder{Interface d'inscription}{fig:ui-register}{6cm}

\figplaceholder{Interface de connexion}{fig:ui-login}{6cm}

\subsection{Gestion du profil}
L'interface de profil (figure~\ref{fig:ui-profile}) permet à l'utilisateur
de consulter et de modifier ses informations personnelles.

\figplaceholder{Interface de gestion du profil}{fig:ui-profile}{6cm}

\subsection{Gestion de l'équipe et des rôles}
L'interface de gestion de l'équipe (figure~\ref{fig:ui-team}) permet à
l'administrateur d'inviter de nouveaux membres, de consulter la liste des
collaborateurs et de leur attribuer un rôle. Trois rôles ont été définis~:
\textbf{administrateur} (droits étendus sur l'entreprise),
\textbf{éditeur} (création et édition de modèles) et \textbf{marketing}
(éditeur habilité à maintenir la galerie de modèles prédéfinis de
l'entreprise). Un rôle machine à machine complète ce dispositif pour les
clients applicatifs authentifiés par clé d'API.

\figplaceholder{Interfaces d'invitation et de gestion de l'équipe}{fig:ui-team}{6cm}

\subsection{Espace de super-administration}
Un espace de super-administration, réservé à l'éditeur de la plateforme,
offre une vue transverse sur l'ensemble des entreprises et de leurs
abonnements. Il permet notamment d'attribuer un plan à une entreprise. Cet
espace, dont les fonctions liées à la facturation sont approfondies au
chapitre~\ref{chap:sprint4}, est présenté à la figure~\ref{fig:ui-superadmin}.

\figplaceholder{Tableau de bord de super-administration}{fig:ui-superadmin}{6cm}

\section*{Conclusion}
Ce sprint a établi le socle de la plateforme : un service d'authentification
robuste, un modèle multi-entreprises garantissant l'isolation des données
par \texttt{tenant\_id}, et une gestion complète des utilisateurs et des
rôles. Ces fondations conditionnent l'ensemble des fonctionnalités
ultérieures. Le sprint suivant s'appuie sur cet acquis pour construire le
c\oe{}ur métier de la plateforme : l'éditeur de modèles multi-canal.

% ---------------------------------------------------------------------
%  CHAPITRE 4
% ---------------------------------------------------------------------
\chapter{Sprint~2 : Éditeur de modèles multi-canal}
\label{chap:sprint2}

\section*{Introduction}
Ce deuxième sprint constitue le c\oe{}ur métier de la plateforme. Il met en
\oe{}uvre l'éditeur de modèles et, en particulier, l'éditeur d'e-mail par
glisser-déposer qui en est le produit phare, ainsi que la prise en charge
des trois canaux (e-mail, facture, contrat) et la bibliothèque de
modèles. Après avoir énoncé les objectifs du sprint, nous détaillons les
cas d'utilisation, les diagrammes de séquence et la réalisation.

\section{Objectifs du sprint}
Ce sprint couvre les epics~4 à~6 du backlog. Le
tableau~\ref{tab:backlog-s2} en présente les user stories.

\begin{table}[H]
\centering
\caption{Backlog du sprint~2}
\label{tab:backlog-s2}
\small
\begin{tabularx}{\linewidth}{clX}
\toprule
\textbf{ID} & \textbf{User story} & \textbf{Description} \\
\midrule
4.1 & Éditer un e-mail & En tant qu'éditeur, je veux composer un e-mail en glissant-déposant des blocs sur un canevas. \\
4.2 & Prévisualiser et coder & En tant qu'éditeur, je veux prévisualiser le rendu en direct et accéder à l'éditeur de code. \\
4.3 & Exporter & En tant qu'éditeur, je veux exporter mon modèle en HTML responsive et en PDF. \\
5.1 & Facture & En tant qu'éditeur, je veux construire une facture structurée avec variables. \\
5.2 & Contrat & En tant qu'éditeur, je veux rédiger un contrat en texte riche paginé. \\
6.1 & Favoris et duplication & En tant qu'éditeur, je veux mettre des modèles en favoris et les dupliquer. \\
6.2 & Galerie prédéfinie & En tant que membre marketing, je veux constituer une galerie de modèles prédéfinis par catégorie. \\
\bottomrule
\end{tabularx}
\end{table}

\section{Raffinement et description des cas d'utilisation}

\subsection{Cas d'utilisation « Gérer les modèles »}
La figure~\ref{fig:uc-templates} regroupe les cas d'utilisation liés à la
gestion des modèles : création, édition, prévisualisation, export,
duplication, mise en favori et gestion de la galerie prédéfinie.

\figplaceholder{Cas d'utilisation « Gérer les modèles »}{fig:uc-templates}{7.5cm}

Le tableau~\ref{tab:uc-create-template} décrit le scénario « Créer un
modèle d'e-mail ».

\begin{table}[H]
\centering
\caption{Description textuelle du cas d'utilisation « Créer un modèle d'e-mail »}
\label{tab:uc-create-template}
\small
\begin{tabularx}{\linewidth}{lX}
\toprule
\textbf{Cas d'utilisation} & Créer un modèle d'e-mail \\
\textbf{Acteur} & Éditeur \\
\textbf{Préconditions} & L'utilisateur est connecté et son plan autorise la création. \\
\textbf{Postconditions} & Un modèle d'e-mail est créé et enregistré pour son entreprise. \\
\midrule
\textbf{Scénario nominal} &
1.~L'éditeur choisit le canal « e-mail » et ouvre l'éditeur.\newline
2.~Il glisse-dépose des blocs sur le canevas et ajuste leur contenu et leur style.\newline
3.~Il prévisualise le rendu en direct.\newline
4.~Il enregistre le modèle.\newline
5.~Le système valide et persiste le modèle rattaché à son entreprise. \\
\midrule
\textbf{Scénario alternatif} &
5.a.~Si le quota du plan est atteint, le système bloque la création et propose une montée en gamme. \\
\bottomrule
\end{tabularx}
\end{table}

Le tableau~\ref{tab:uc-export} décrit le scénario « Exporter un modèle ».

\begin{table}[H]
\centering
\caption{Description textuelle du cas d'utilisation « Exporter un modèle »}
\label{tab:uc-export}
\small
\begin{tabularx}{\linewidth}{lX}
\toprule
\textbf{Cas d'utilisation} & Exporter un modèle \\
\textbf{Acteur} & Éditeur \\
\textbf{Préconditions} & Un modèle existe et est ouvert. \\
\textbf{Postconditions} & Un fichier HTML ou PDF est généré et téléchargé. \\
\midrule
\textbf{Scénario nominal} &
1.~L'éditeur demande l'export dans le format souhaité.\newline
2.~Le système compile le contenu (MJML vers HTML pour l'e-mail) et substitue les variables.\newline
3.~Le système produit le fichier et le met à disposition de l'utilisateur. \\
\midrule
\textbf{Scénario alternatif} &
2.a.~En cas d'erreur de compilation, le système en informe l'utilisateur. \\
\bottomrule
\end{tabularx}
\end{table}

\section{Diagrammes de séquence}

\subsection{Séquence « Créer un modèle »}
La figure~\ref{fig:seq-create-template} décrit les échanges entre l'éditeur,
la passerelle API et le service de modèles lors de la création d'un modèle,
depuis la saisie jusqu'à la persistance en base.

\figplaceholder{Diagramme de séquence « Créer un modèle »}{fig:seq-create-template}{7cm}

\subsection{Séquence « Exporter un modèle »}
La figure~\ref{fig:seq-export} illustre le processus d'export : compilation
du contenu, substitution des variables et génération du fichier.

\figplaceholder{Diagramme de séquence « Exporter un modèle »}{fig:seq-export}{7cm}

\section{Réalisation}

\subsection{Service de modèles et modèle de données par blocs}
Le service de modèles, développé avec NestJS selon le patron CQRS, expose en
gRPC l'ensemble des opérations du cycle de vie d'un modèle : création,
modification, suppression, consultation, listage, rendu, duplication,
gestion des favoris et récupération des modèles populaires. Les données sont
persistées dans la base \texttt{templates\_db}~; les médias (images) sont
stockés dans MinIO.

Chaque modèle est caractérisé par un \textbf{type de canal} — e-mail (1),
facture (2) ou contrat (3) — et par un contenu dont la
forme dépend du canal. Pour l'e-mail, le contenu suit un
\textbf{modèle de blocs} structuré : un modèle est composé de rangées
(\emph{rows}), chaque rangée pouvant comporter jusqu'à quatre colonnes, et
chaque colonne contenant des blocs. Onze types de blocs sont disponibles :
titre, texte, image, vidéo, bouton, séparateur, tableau, signature, réseaux
sociaux, menu et liste à icônes. Ce modèle de données présente une propriété
essentielle : il est \textbf{identique} à celui que produit l'assistant
d'IA (chapitre~\ref{chap:sprint3}), de sorte que l'édition manuelle et
l'édition assistée deviennent parfaitement interchangeables.

\subsection{Liste des modèles et bibliothèque}
La figure~\ref{fig:ui-templates-list} présente la liste des modèles de
l'entreprise, avec aperçu, mise en favori, duplication et suppression. Une
page dédiée regroupe les favoris, et une galerie expose les modèles
prédéfinis maintenus par les membres marketing, organisés par catégorie.

\figplaceholder{Liste des modèles}{fig:ui-templates-list}{6cm}

\subsection{Éditeur d'e-mail visuel}
L'éditeur d'e-mail (figure~\ref{fig:ui-email-editor}) est le composant phare
de la plateforme. Il s'organise autour d'un canevas central sur lequel
l'utilisateur glisse-dépose des blocs, d'un panneau de blocs à gauche et
d'un panneau de style à droite permettant d'ajuster les propriétés (couleur,
espacement, alignement, largeurs de colonnes, styles de section). Le
glisser-déposer s'appuie sur la bibliothèque dnd-kit. Un aperçu de l'e-mail
est rendu en direct, et un éditeur de code (Monaco) est mis à disposition
des utilisateurs avancés. À l'export, le modèle est compilé en HTML
responsive au moyen de MJML, ou en PDF.

\figplaceholder{Éditeur d'e-mail visuel}{fig:ui-email-editor}{7cm}

\subsection{Éditeurs de facture et de contrat}
Les canaux transactionnel et documentaire disposent chacun d'un éditeur
adapté. L'éditeur de \textbf{facture} propose un constructeur structuré avec
variables de personnalisation et génération PDF. L'éditeur de
\textbf{contrat} repose sur un composant de texte riche (Tiptap) avec
pagination et export PDF. La figure~\ref{fig:ui-facture-contrat} en présente
un aperçu.

\figplaceholder{Éditeurs de facture et de contrat}{fig:ui-facture-contrat}{6cm}

\section*{Conclusion}
Ce sprint a livré le c\oe{}ur fonctionnel de la plateforme : un éditeur
d'e-mail visuel complet reposant sur un modèle de blocs structuré, la prise
en charge des trois canaux de communication, l'export HTML et PDF, ainsi que
la bibliothèque de modèles (favoris, duplication, galerie prédéfinie). Le
caractère structuré et partagé du modèle de blocs prépare directement le
sprint suivant, consacré à l'assistant d'intelligence artificielle, qui
produira et modifiera ces mêmes blocs.

% ---------------------------------------------------------------------
%  CHAPITRE 5
% ---------------------------------------------------------------------
\chapter{Sprint~3 : Assistant d'intelligence artificielle agentique}
\label{chap:sprint3}

\section*{Introduction}
Ce sprint porte la principale valeur ajoutée du projet : un assistant
conversationnel capable de générer et de modifier des modèles d'e-mail en
manipulant de véritables blocs éditables, et même de reconstruire une
campagne à partir d'une affiche. Nous présentons d'abord les objectifs du
sprint, puis la conception de l'assistant — modèles employés, architecture
de génération, protocole d'exposition des outils et chaîne de traitement de
l'image —, avant de décrire les diagrammes de séquence et la réalisation.

\section{Objectifs du sprint}
Ce sprint couvre les epics~7 à~10 du backlog. Le
tableau~\ref{tab:backlog-s3} en détaille les user stories.

\begin{table}[H]
\centering
\caption{Backlog du sprint~3}
\label{tab:backlog-s3}
\small
\begin{tabularx}{\linewidth}{clX}
\toprule
\textbf{ID} & \textbf{User story} & \textbf{Description} \\
\midrule
7.1 & Générer par consigne & En tant qu'éditeur, je veux qu'un e-mail complet et à ma marque soit généré à partir d'une consigne en langage naturel. \\
8.1 & Éditer en conversation & En tant qu'éditeur, je veux modifier mon e-mail par dialogue (thème, image, disposition, ajout ou déplacement de blocs). \\
8.2 & Éditer par sélection & En tant qu'éditeur, je veux sélectionner un élément et lui appliquer une consigne ciblée. \\
9.1 & Affiche vers e-mail & En tant qu'éditeur, je veux importer une affiche et obtenir un e-mail éditable reprenant la campagne. \\
10.1 & Exposition des outils (MCP) & En tant que système, je veux exposer les outils de construction via un protocole standardisé afin que le modèle les découvre dynamiquement. \\
10.2 & Modèle auto-hébergé & En tant qu'entreprise, je veux que l'inférence s'exécute sur un serveur privé afin de préserver la confidentialité des données. \\
\bottomrule
\end{tabularx}
\end{table}

\section{Conception de l'assistant}

\subsection{Le modèle de blocs comme contrat d'outils}
Le principe fondateur de l'assistant est le suivant : au lieu de produire un
bloc de HTML monolithique, le modèle de langage \textbf{appelle des outils}
de construction qui renvoient exactement le même format de blocs
(\texttt{BlockData}) que l'éditeur manuel décrit au
chapitre~\ref{chap:sprint2}. Le HTML n'est généré qu'au moment de l'export.
Cette approche, dite \emph{agentique}, présente deux avantages décisifs :
les modifications visuelles et les modifications par IA opèrent sur la même
structure de données (elles sont donc interchangeables), et le résultat est
immédiatement et finement éditable par l'utilisateur.

Un ensemble de \textbf{32 outils} a été défini : un outil par type de bloc,
complété par des outils transversaux (définition du thème, ouverture de
section, carte, en-tête sur image, barre de couleurs de marque, changement
de disposition, déplacement, édition et suppression de blocs).

\subsection{Modèles et infrastructure}
L'inférence repose sur un \textbf{serveur de modèle de langage
auto-hébergé} exploitant le moteur haute performance vLLM, exposé via une
API compatible OpenAI et protégé par authentification. Deux modèles à poids
ouverts y sont servis : \textbf{Gemma} (modèle par défaut) et un modèle de
repli, tous deux dotés d'un large contexte. Ce choix d'auto-hébergement est
central : il garantit qu'aucune donnée de l'entreprise n'est transmise à une
API tierce, répondant ainsi à l'exigence de souveraineté formulée dans les
besoins non fonctionnels.

\subsection{Architecture de génération : planificateur et exécuteur}
Le chemin principal de génération d'un e-mail neuf repose sur un
\textbf{planificateur déterministe}. Plutôt que de laisser le modèle
enchaîner librement des dizaines d'appels d'outils — approche fragile avec
un modèle de taille modérée —, une \textbf{unique requête contrainte par un
schéma} produit une spécification de design (\emph{DesignSpec}) : système de
design, ambiance, palette de couleurs et liste ordonnée de sections. Cette
spécification est ensuite \textbf{exécutée par le code} pour construire les
blocs de manière déterministe. Il en résulte une génération rapide, sans
duplication et avec un pied de page unique. Lorsque le planificateur échoue
ou renvoie un résultat trop pauvre, une \textbf{voie agentique de repli}
prend le relais, laissant le modèle appeler directement les outils. La
figure~\ref{fig:archi-ia} synthétise cette architecture.

\figplaceholder{Architecture de l'assistant d'intelligence artificielle}{fig:archi-ia}{7.5cm}

\subsection{Exposition des outils via le protocole MCP}
Pour les tours pilotés par le modèle (édition et repli agentique), les
outils de construction sont exposés au moyen du \textbf{Model Context
Protocol} (MCP). Un point d'accès dédié agit comme serveur MCP : le modèle y
\textbf{découvre dynamiquement} les 32 outils disponibles plutôt que de se
les voir énumérer dans l'invite. Cette décorrélation entre la définition des
outils et l'invite constitue une architecture agentique moderne et
standardisée. Le serveur MCP réutilise exactement les mêmes définitions
d'outils que l'éditeur, garantissant une source de vérité unique.

\subsection{Chaîne « affiche vers e-mail »}
La fonction d'import d'affiche traite une image en \textbf{trois étapes}.
\begin{enumerate}[leftmargin=1.4cm]
  \item \textbf{Lecture par vision.} L'image est analysée par le modèle
  multimodal, qui en extrait un rapport structuré : texte exact (OCR),
  offre (prix, cadeaux, forfaits, code promotionnel, appel à l'action),
  couleurs dominantes, ambiance, marque et disposition. Cette étape, la seule
  réellement multimodale, s'appuie sur un décodage contraint par schéma pour
  fiabiliser l'extraction.
  \item \textbf{Rapport vers directive.} Le rapport et la consigne de
  l'utilisateur sont transformés en une directive de génération et en une
  amorce de thème, qui alimentent le générateur de blocs déjà décrit :
  l'image gouverne la marque et le contenu, la consigne oriente la
  structure.
  \item \textbf{Passe critique.} Une passe d'auto-vérification compare
  l'e-mail produit à l'affiche d'origine et signale les écarts (prix erroné,
  élément manquant), corrigés par une passe finale.
\end{enumerate}
Il est essentiel de noter que la plateforme ne colle jamais l'affiche : elle
\textbf{reconstruit} nativement la campagne en blocs éditables. La
figure~\ref{fig:pipeline-image} illustre cette chaîne de traitement.

\figplaceholder{Chaîne de traitement « affiche vers e-mail »}{fig:pipeline-image}{6cm}

\subsection{Intelligence de design et fiabilité}
Plusieurs mécanismes garantissent la qualité et la fiabilité du résultat.
\begin{itemize}[leftmargin=1.4cm]
  \item \textbf{Cinq systèmes de design} (éditorial, audacieux, minimal,
  luxe, corporate) paramètrent l'échelle typographique, les alignements, les
  espacements et le traitement des boutons et des cartes~; le système est
  choisi en fonction du sujet.
  \item Un \textbf{dériveur de palette} et une \textbf{garde de contraste}
  assurent que le texte reste toujours lisible sur son fond, quelle que soit
  la couleur choisie.
  \item Un \textbf{routeur d'intention} classe chaque tour (effacer,
  réécrire, image, éditer, changer de thème) afin d'aiguiller la requête vers
  le traitement — déterministe ou piloté par le modèle — le plus adapté.
  \item Des \textbf{confirmations honnêtes} garantissent que l'assistant
  n'annonce un succès que si le modèle a réellement été modifié, un
  compteur de mutations servant de preuve de changement.
\end{itemize}

\section{Diagrammes de séquence}

\subsection{Séquence « Générer un e-mail par consigne »}
La figure~\ref{fig:seq-generate} décrit le flux de génération : la consigne
est transmise à la route d'IA, le planificateur produit une spécification de
design via le serveur d'inférence, le code exécute cette spécification en
blocs, les images sont résolues, puis le modèle est renvoyé au client et
appliqué au canevas.

\figplaceholder{Diagramme de séquence « Générer un e-mail par consigne »}{fig:seq-generate}{7.5cm}

\subsection{Séquence « Reconstruire depuis une affiche »}
La figure~\ref{fig:seq-image} illustre la chaîne complète de l'import
d'affiche : lecture par vision, transformation en directive, génération des
blocs et passe critique.

\figplaceholder{Diagramme de séquence « Reconstruire depuis une affiche »}{fig:seq-image}{7.5cm}

\section{Réalisation}

\subsection{Panneau de discussion et génération}
L'assistant se présente sous la forme d'un panneau de discussion intégré à
l'éditeur d'e-mail (figure~\ref{fig:ui-ai-chat}). À partir d'une simple
consigne — par exemple « crée un e-mail Black Friday pour notre boutique » —
l'assistant produit en une dizaine de secondes un e-mail complet et à la
marque : en-tête, section d'offre, grille de produits et pied de page.

\figplaceholder{Panneau de discussion IA et e-mail généré}{fig:ui-ai-chat}{6.5cm}

\subsection{Édition conversationnelle et ciblée}
L'utilisateur peut ensuite affiner son e-mail par le dialogue : « passe-le
en thème sombre », « change la première image », « déplace la section des
produits avant l'offre », « ajoute une boîte de prix ». L'assistant vise le
bon bloc et applique la modification. Il est également possible de
\textbf{sélectionner} un bloc ou une section dans le canevas puis de saisir
une consigne : la modification s'applique alors précisément à l'élément
sélectionné. La figure~\ref{fig:ui-ai-edit} illustre l'édition ciblée par
sélection.

\figplaceholder{Édition ciblée par sélection}{fig:ui-ai-edit}{6cm}

\subsection{Import d'une affiche}
La figure~\ref{fig:ui-image-import} présente l'import d'une affiche : après
téléversement, l'assistant affiche l'avancement des étapes (lecture, analyse,
génération) puis produit l'e-mail éditable reconstruit, fidèle à la marque
et à l'offre de l'affiche d'origine.

\figplaceholder{Import d'une affiche et e-mail reconstruit}{fig:ui-image-import}{6.5cm}

\section*{Conclusion}
Ce sprint a doté la plateforme de son principal différenciateur : un
assistant agentique qui génère et modifie de véritables blocs éditables,
reconstruit une campagne à partir d'une affiche grâce à la vision, et
s'appuie sur un modèle de langage auto-hébergé garantissant la souveraineté
des données. L'architecture combine un planificateur déterministe rapide,
une voie agentique de repli et une exposition des outils via le protocole
MCP, le tout renforcé par une intelligence de design et des mécanismes de
fiabilité. Le sprint suivant transforme cette plateforme fonctionnellement
riche en un service commercialisable : abonnements, intégrations et mise en
production.

% ---------------------------------------------------------------------
%  CHAPITRE 6
% ---------------------------------------------------------------------
\chapter{Sprint~4 : Abonnements, intégrations et mise en production}
\label{chap:sprint4}

\section*{Introduction}
Ce dernier sprint transforme une plateforme fonctionnellement complète en un
service exploitable et commercialisable. Il met en place le système
d'abonnements et de facturation, l'application des quotas par plan, la
supervision par le super-administrateur, la couche d'intégration permettant
d'embarquer l'éditeur, le site marketing public, et enfin le déploiement en
production. Nous présentons les objectifs du sprint, les cas d'utilisation,
les diagrammes de séquence, puis la réalisation.

\section{Objectifs du sprint}
Ce sprint couvre les epics~11 à~15 du backlog. Le
tableau~\ref{tab:backlog-s4} en présente les user stories.

\begin{table}[H]
\centering
\caption{Backlog du sprint~4}
\label{tab:backlog-s4}
\small
\begin{tabularx}{\linewidth}{clX}
\toprule
\textbf{ID} & \textbf{User story} & \textbf{Description} \\
\midrule
11.1 & Souscrire un abonnement & En tant qu'administrateur, je veux souscrire un plan payant via un paiement en ligne sécurisé. \\
11.2 & Gérer l'abonnement & En tant qu'administrateur, je veux changer, annuler ou réactiver mon abonnement. \\
11.3 & Quotas par plan & En tant que système, je veux appliquer les limites propres à chaque plan. \\
12.1 & Attribuer les plans & En tant que super-administrateur, je veux superviser les entreprises et leur attribuer un plan. \\
13.1 & Clés d'API & En tant qu'administrateur, je veux générer des clés d'intégration. \\
13.2 & Embarquer l'éditeur & En tant qu'outil externe, je veux envoyer mes utilisateurs dans l'éditeur sans seconde authentification. \\
14.1 & Site marketing & En tant que visiteur, je veux consulter un site public présentant la plateforme et ses offres. \\
15.1 & Mise en production & En tant qu'équipe, je veux déployer la plateforme de façon sécurisée et reproductible. \\
\bottomrule
\end{tabularx}
\end{table}

\section{Raffinement et description des cas d'utilisation}

\subsection{Cas d'utilisation « Gérer son abonnement »}
La figure~\ref{fig:uc-billing} regroupe les cas d'utilisation liés à la
facturation : souscription, changement de plan, annulation, réactivation, et
attribution de plan par le super-administrateur.

\figplaceholder{Cas d'utilisation « Gérer son abonnement »}{fig:uc-billing}{7cm}

Le tableau~\ref{tab:uc-subscribe} décrit le scénario « Souscrire un
abonnement ».

\begin{table}[H]
\centering
\caption{Description textuelle du cas d'utilisation « Souscrire un abonnement »}
\label{tab:uc-subscribe}
\small
\begin{tabularx}{\linewidth}{lX}
\toprule
\textbf{Cas d'utilisation} & Souscrire un abonnement \\
\textbf{Acteur} & Administrateur d'entreprise \\
\textbf{Préconditions} & L'administrateur est connecté~; son entreprise est sur un plan inférieur. \\
\textbf{Postconditions} & L'entreprise est passée au plan payant choisi. \\
\midrule
\textbf{Scénario nominal} &
1.~L'administrateur choisit un plan et une périodicité (mensuelle ou annuelle).\newline
2.~Le système ouvre le paiement intégré et présente le détail hors taxes et TTC.\newline
3.~L'administrateur renseigne son moyen de paiement et confirme.\newline
4.~Le prestataire de paiement valide la transaction.\newline
5.~Le système met à jour le plan de l'entreprise et débloque les fonctionnalités correspondantes. \\
\midrule
\textbf{Scénario alternatif} &
4.a.~En cas d'échec du paiement, le système en informe l'administrateur, qui peut réessayer. \\
\bottomrule
\end{tabularx}
\end{table}

\section{Diagrammes de séquence}

\subsection{Séquence « Souscrire un abonnement »}
La figure~\ref{fig:seq-subscribe} décrit l'enchaînement entre le client, la
passerelle API, le prestataire de paiement et le service
d'authentification lors d'une souscription, jusqu'à la mise à jour du plan
de l'entreprise.

\figplaceholder{Diagramme de séquence « Souscrire un abonnement »}{fig:seq-subscribe}{7cm}

\subsection{Séquence « Embarquer l'éditeur »}
La figure~\ref{fig:seq-embed} illustre le flux d'intégration : un outil
externe redirige son utilisateur vers l'éditeur avec une authentification
déléguée, celui-ci construit un modèle, puis est redirigé vers l'outil
d'origine.

\figplaceholder{Diagramme de séquence « Embarquer l'éditeur »}{fig:seq-embed}{7cm}

\section{Réalisation}

\subsection{Plans et grille tarifaire}
Quatre plans ont été définis : \textbf{Free} (gratuit), \textbf{Pro},
\textbf{Pro Organisation}, ainsi qu'un plan \textbf{interne} non
commercialisable réservé aux entités du groupe. Le
tableau~\ref{tab:plans} récapitule leurs caractéristiques.

\begin{table}[H]
\centering
\caption{Plans d'abonnement de la plateforme}
\label{tab:plans}
\small
\begin{tabularx}{\linewidth}{lcX}
\toprule
\textbf{Plan} & \textbf{Tarif (HT/mois)} & \textbf{Fonctionnalités} \\
\midrule
Free & Gratuit & Canal e-mail uniquement, 1 modèle, 1 interaction IA. \\
Pro & 25~€ & Tous les canaux, IA illimitée (hors invitations et clés d'API). \\
Pro Organisation & 55~€ & Pro + gestion d'équipe + accès API / intégrations. \\
Interne & — & Illimité~; attribué par le super-administrateur aux entités du groupe. \\
\bottomrule
\end{tabularx}
\end{table}

La grille tarifaire publique (figure~\ref{fig:ui-pricing}) propose une
bascule entre facturation mensuelle et annuelle. La \textbf{TVA} française
(20~\%) est correctement appliquée, le détail hors taxes et toutes taxes
comprises étant présenté au moment du paiement.

\figplaceholder{Page de tarification}{fig:ui-pricing}{6cm}

\subsection{Paiement et facturation}
Le paiement s'appuie sur le prestataire Stripe, au moyen d'un
\textbf{Checkout intégré} : le règlement s'effectue sur la page de la
plateforme, sans redirection externe (figure~\ref{fig:ui-checkout}). Les
abonnements sont récurrents et peuvent être mensuels (flexibles) ou annuels
(engagement de douze mois, à tarif réduit). Un onglet de facturation permet à
l'administrateur de consulter son plan, de le changer, de l'annuler dans le
respect de l'engagement, ou de le réactiver. Les événements de paiement sont
répercutés sur le plan de l'entreprise, tant par confirmation directe que par
\emph{webhook} de sécurité.

\figplaceholder{Paiement intégré et onglet de facturation}{fig:ui-checkout}{6cm}

\subsection{Application des quotas par plan}
Les limites de chaque plan sont appliquées côté passerelle API à partir
d'une source unique de vérité. Le plan gratuit est restreint au canal
e-mail, à un modèle et à une interaction IA, ce décompte étant monotone (il
n'est jamais décrémenté, de sorte qu'une suppression puis recréation ne
permet pas de contourner la limite). Lorsqu'une limite est atteinte, une
\textbf{fenêtre de montée en gamme} invite l'utilisateur à passer à un plan
supérieur (figure~\ref{fig:ui-upgrade}).

\figplaceholder{Fenêtre de montée en gamme}{fig:ui-upgrade}{5cm}

\subsection{Supervision par le super-administrateur}
L'espace de super-administration a été enrichi d'une vue des commandes et
d'indicateurs clés (revenu récurrent, entreprises actives, impayés) obtenus
par enrichissement des données de facturation. Le super-administrateur peut y
attribuer un plan à toute entreprise, y compris le plan interne. La
figure~\ref{fig:ui-superadmin-billing} présente cet espace.

\figplaceholder{Supervision des abonnements}{fig:ui-superadmin-billing}{6cm}

\subsection{Intégrations et embarquement}
Un administrateur peut générer des \textbf{clés d'API} depuis le panneau
« Intégrations » des paramètres. Un \textbf{flux d'embarquement} permet à un
outil externe (par exemple un CRM) d'envoyer ses utilisateurs dans l'éditeur
avec une authentification déléguée — sur le principe d'un Checkout délégué,
sans seconde connexion —, de leur faire construire un modèle, puis de les
rediriger vers l'outil d'origine, le tout cloisonné à l'entreprise
concernée.

\subsection{Site marketing public}
Un site marketing public a été développé pour présenter la plateforme
(accueil, fonctionnalités, tarifs, à propos, contact, démonstration). Sa page
de tarification est reliée aux plans réels de la plateforme. La
figure~\ref{fig:ui-marketing} en présente la page d'accueil.

\figplaceholder{Site marketing public}{fig:ui-marketing}{6cm}

\subsection{Mise en production}
L'ensemble des services est \textbf{conteneurisé} avec Docker et orchestré
par Docker Compose, deux piles distinctes étant maintenues pour le
développement et la production. La plateforme est \textbf{déployée en
production} sur un serveur du groupe. Son exposition publique repose sur un
\textbf{tunnel Cloudflare}, sécurisé et sortant, qui ne nécessite l'ouverture
d'aucun port entrant et confie la terminaison HTTPS à Cloudflare. Plusieurs
noms d'hôte publics desservent respectivement le frontend, la passerelle API,
le stockage MinIO et les services d'IA. Le code est hébergé sur GitLab, les
secrets étant tenus hors du dépôt. La figure~\ref{fig:archi-deploy} présente
l'architecture de déploiement.

\figplaceholder{Architecture de déploiement}{fig:archi-deploy}{6.5cm}

\section*{Conclusion}
Ce sprint a doté la plateforme des attributs d'un véritable service SaaS de
production : un système d'abonnements Stripe avec TVA et gestion complète du
cycle de vie, l'application rigoureuse des quotas par plan, la supervision
par le super-administrateur, une couche d'intégration et d'embarquement, un
site marketing public, et un déploiement sécurisé via Docker et Cloudflare.
La plateforme est ainsi passée de la preuve de concept fonctionnelle au
produit exploitable en conditions réelles.

% ---------------------------------------------------------------------
%  CONCLUSION GÉNÉRALE
% ---------------------------------------------------------------------
\chapter*{Conclusion générale}
\addcontentsline{toc}{chapter}{Conclusion générale}
\markboth{Conclusion générale}{Conclusion générale}

Réalisé au sein de France Téléphone, filiale du groupe Winvest Capital, ce
projet de fin d'études avait pour ambition de concevoir et de développer,
de bout en bout, une plateforme SaaS multi-entreprises de création de
modèles de communication assistée par intelligence artificielle. Au terme
de ce travail, l'objectif fixé a été atteint : \textbf{Winaity Template
Builder} est une plateforme fonctionnelle et déployée en production.

\paragraph{Récapitulation de la démarche.}
La conduite du projet a suivi la méthodologie agile Scrum, structurée en
quatre sprints. Après une phase d'analyse des besoins et de conception —
identification des acteurs, expression des besoins, modélisation UML et
définition de l'architecture microservices —, nous avons successivement mis
en \oe{}uvre l'authentification et le modèle multi-entreprises (sprint~1),
l'éditeur de modèles multi-canal (sprint~2), l'assistant d'intelligence
artificielle agentique (sprint~3), puis les abonnements, les intégrations et
la mise en production (sprint~4).

\paragraph{Résultats obtenus.}
La plateforme répond à la problématique posée en introduction. Elle permet à
une équipe métier, sans compétence technique, de produire rapidement des
supports sur trois canaux (e-mail, facture, contrat) au moyen d'un
éditeur visuel. Son assistant d'IA génère et modifie de véritables blocs
éditables et sait reconstruire une campagne à partir d'une simple affiche.
L'isolation multi-entreprises garantit la confidentialité des données de
chaque société, tandis que l'auto-hébergement du modèle de langage assure la
souveraineté des données — deux exigences qu'aucune des solutions du marché
étudiées ne réunissait. La plateforme intègre enfin la gestion des rôles, un
système d'abonnements Stripe conforme sur le plan fiscal, une couche
d'intégration, et un déploiement sécurisé.

\paragraph{Difficultés rencontrées.}
La principale difficulté a résidé dans la fiabilisation de l'assistant
d'intelligence artificielle. Un modèle de langage de taille modérée, servi
localement, se révèle peu fiable lorsqu'on le laisse enchaîner librement de
nombreux appels d'outils : duplications, omissions ou actions inopérantes.
Cette difficulté a été surmontée par une architecture hybride associant un
planificateur déterministe contraint par schéma, une voie agentique de
repli, un routeur d'intention et des mécanismes de fiabilité (garde de
contraste, confirmations honnêtes, filtrage des sorties du modèle).
L'exécution multimodale (lecture d'affiche) a également imposé d'isoler
l'étape de vision et de la fiabiliser par décodage contraint. Sur le plan de
l'ingénierie, la coordination d'une architecture microservices et sa mise en
production ont nécessité une attention particulière à la configuration, aux
secrets et à la reproductibilité des environnements.

\paragraph{Apports du stage.}
Ce projet a été l'occasion d'acquérir et d'approfondir de nombreuses
compétences. Sur le plan technique : la maîtrise d'une architecture
microservices moderne (NestJS, gRPC, CQRS, TypeORM, PostgreSQL, MinIO), du
développement frontend avec Next.js et React, de l'intégration d'un système
de paiement, et surtout de l'ingénierie appliquée aux grands modèles de
langage — génération agentique par outils, planification contrainte par
schéma, vision par ordinateur et protocole MCP, sur une infrastructure
d'inférence auto-hébergée. Sur le plan professionnel et personnel, ce stage
a permis d'évoluer au sein d'une équipe, de dialoguer avec des interlocuteurs
métier, de conduire un projet de son analyse à sa mise en production, et de
prendre des décisions d'architecture argumentées.

\paragraph{Perspectives.}
Plusieurs pistes d'évolution ont été identifiées. Sur le plan collaboratif,
l'introduction de fonctionnalités temps réel (curseurs et synchronisation des
brouillons) enrichirait le travail à plusieurs. Sur le plan de l'IA, le
recours à un modèle plus puissant pour la seule étape de planification
améliorerait la fidélité des mises en page complexes. Sur le
plan de l'écosystème, l'approfondissement du flux d'intégration permettrait
de connecter la plateforme à des outils tiers tels que des CRM, et le site
marketing pourrait être enrichi (contenu, référencement, démonstrations). Ces
perspectives confirment le potentiel de la plateforme et ouvrent la voie à
son développement continu.

% ---------------------------------------------------------------------
%  ANNEXES
% ---------------------------------------------------------------------
\appendix
\cleardoublepage

\chapter{Structure de données d'un bloc (\emph{BlockData})}
\label{annexe:blockdata}
Le modèle de données ci-dessous est partagé entre l'éditeur manuel et
l'assistant d'intelligence artificielle : c'est la même structure qui est
produite dans les deux cas. Un modèle d'e-mail est un ensemble de rangées~;
chaque rangée contient des colonnes~; chaque colonne contient des blocs.

\begin{lstlisting}
// Bloc atomique : meme structure produite par l'editeur et par l'IA
interface BlockData {
  id: string;
  type: 'heading' | 'text' | 'image' | 'video' | 'button'
      | 'divider' | 'table' | 'signature' | 'social'
      | 'menu' | 'icon-list';
  content: Record<string, unknown>;   // texte, src, items, ...
  styles:  Record<string, string>;    // couleur, padding, alignement, ...
}

interface Column   { id: string; width: string; blocks: BlockData[]; styles?: object; }
interface Row      { id: string; layout: string; columns: Column[]; styles: object; }
interface Template { rows: Row[]; globalStyles: object; }
\end{lstlisting}

\chapter{Extrait de la configuration Docker Compose (production)}
\label{annexe:compose}
Extrait simplifié de l'orchestration des services en production. Les
secrets ne figurent pas dans le fichier~: ils sont fournis par des
variables d'environnement.

\begin{lstlisting}
services:
  frontend:          # Next.js (interface + routes IA)
    build: ./frontend
    environment: [ NEXT_PUBLIC_API_URL, AI_BASE_URL ]
  api-gateway:       # unique surface REST publique -> gRPC
    build: ./api-gateway
    environment: [ JWT_SECRET, FRONTEND_ORIGIN, STRIPE_SECRET_KEY ]
    depends_on: [ auth-service, template-service ]
  auth-service:      # comptes, tenants, roles, abonnements (gRPC)
    build: ./auth-service
    depends_on: [ postgres ]
  template-service:  # CRUD + rendu des modeles (gRPC, CQRS)
    build: ./template-service
    depends_on: [ postgres, minio ]
  postgres:          # bases auth_db + templates_db
    image: postgres:18
  minio:             # stockage objet (medias)
    image: minio/minio
  cloudflared:       # tunnel Cloudflare (exposition HTTPS sortante)
    image: cloudflare/cloudflared
\end{lstlisting}

% ---------------------------------------------------------------------
%  BIBLIOGRAPHIE / NETOGRAPHIE  (manuelle — pas de biber)
% ---------------------------------------------------------------------
\cleardoublepage
\addcontentsline{toc}{chapter}{Bibliographie / Netographie}
\begin{thebibliography}{99}

\bibitem{nextjs} Vercel, \emph{Next.js Documentation}. \url{https://nextjs.org/docs} (consulté en 2026).
\bibitem{react} Meta Open Source, \emph{React — The library for web and native user interfaces}. \url{https://react.dev/} (consulté en 2026).
\bibitem{nestjs} NestJS, \emph{A progressive Node.js framework}. \url{https://docs.nestjs.com/} (consulté en 2026).
\bibitem{grpc} Cloud Native Computing Foundation, \emph{gRPC — A high performance, open source universal RPC framework}. \url{https://grpc.io/docs/} (consulté en 2026).
\bibitem{postgresql} The PostgreSQL Global Development Group, \emph{PostgreSQL Documentation}. \url{https://www.postgresql.org/docs/} (consulté en 2026).
\bibitem{typeorm} TypeORM, \emph{Amazing ORM for TypeScript and JavaScript}. \url{https://typeorm.io/} (consulté en 2026).
\bibitem{minio} MinIO Inc., \emph{High Performance Object Storage}. \url{https://min.io/docs/} (consulté en 2026).
\bibitem{mjml} Mailjet, \emph{MJML — The only framework that makes responsive email easy}. \url{https://mjml.io/} (consulté en 2026).
\bibitem{vercelaisdk} Vercel, \emph{AI SDK — The AI Toolkit for TypeScript}. \url{https://ai-sdk.dev/} (consulté en 2026).
\bibitem{mcp} Anthropic, \emph{Model Context Protocol (MCP) Specification}. \url{https://modelcontextprotocol.io/} (consulté en 2026).
\bibitem{vllm} vLLM Project, \emph{Easy, fast, and cheap LLM serving}. \url{https://docs.vllm.ai/} (consulté en 2026).
\bibitem{gemma} Google DeepMind, \emph{Gemma open models}. \url{https://ai.google.dev/gemma} (consulté en 2026).
\bibitem{stripe} Stripe, \emph{Documentation — Payments, Billing, Checkout}. \url{https://docs.stripe.com/} (consulté en 2026).
\bibitem{docker} Docker Inc., \emph{Docker Documentation}. \url{https://docs.docker.com/} (consulté en 2026).
\bibitem{cloudflaretunnel} Cloudflare, \emph{Cloudflare Tunnel Documentation}. \url{https://developers.cloudflare.com/cloudflare-one/} (consulté en 2026).
\bibitem{scrumguide} K. Schwaber et J. Sutherland, \emph{The Scrum Guide}, 2020. \url{https://scrumguides.org/} (consulté en 2026).
\bibitem{tailwind} Tailwind Labs, \emph{Tailwind CSS Documentation}. \url{https://tailwindcss.com/docs} (consulté en 2026).
\bibitem{dndkit} dnd kit, \emph{A modern drag and drop toolkit for React}. \url{https://dndkit.com/} (consulté en 2026).

\end{thebibliography}

% ---------------------------------------------------------------------
%  RÉSUMÉ / ABSTRACT
% ---------------------------------------------------------------------
\cleardoublepage
\thispagestyle{empty}

\section*{Résumé}
Ce rapport présente la conception et le développement de
\textbf{Winaity Template Builder}, une plateforme SaaS multi-entreprises
permettant de créer, gérer et générer des modèles de communication
professionnels sur trois canaux — e-mail, facture et contrat — à travers un
éditeur visuel par glisser-déposer enrichi d'un
assistant d'intelligence artificielle \emph{agentique}. L'assistant génère
et modifie de véritables blocs éditables (et non du HTML jetable) au moyen
d'un ensemble d'outils, et peut reconstruire une campagne à partir d'une
simple affiche marketing grâce à la vision par ordinateur. La plateforme
repose sur une architecture microservices (passerelle API, service
d'authentification et service de modèles communiquant en gRPC), une base
PostgreSQL, un stockage objet MinIO, et un serveur de modèle de langage
auto-hébergé (vLLM / Gemma) garantissant la souveraineté des données. Elle
intègre l'isolation multi-locataire, la gestion des rôles, un système
d'abonnements Stripe avec TVA, ainsi qu'un déploiement en production. Le
projet a été conduit selon la méthodologie agile Scrum au sein de
France Téléphone (groupe Winvest Capital).

\smallskip
\noindent\textbf{Mots-clés :} Intelligence Artificielle agentique, SaaS
multi-entreprises, éditeur de modèles, e-mail, MJML, microservices, gRPC,
Next.js, NestJS, PostgreSQL, vLLM, MCP, Stripe.

\vspace{1cm}

\section*{Abstract}
This report presents the design and development of
\textbf{Winaity Template Builder}, a multi-tenant SaaS platform for
creating, managing and generating professional communication templates
across three channels — email, invoice and contract — through a
drag-and-drop visual editor powered by an \emph{agentic}
artificial-intelligence assistant. The assistant generates and edits real
editable blocks (not throwaway HTML) through a set of tools, and can rebuild
a marketing campaign from a single poster image using computer vision. The
platform is built on a microservices architecture (API gateway,
authentication service and template service communicating over gRPC), a
PostgreSQL database, MinIO object storage, and a self-hosted large language
model server (vLLM / Gemma) ensuring data sovereignty. It provides
multi-tenant isolation, role-based access control, a Stripe subscription
system with VAT, and a production deployment. The project was carried out
using the agile Scrum methodology at France Téléphone (Winvest Capital
group).

\smallskip
\noindent\textbf{Keywords:} Agentic Artificial Intelligence, Multi-tenant
SaaS, Template builder, Email, MJML, Microservices, gRPC, Next.js, NestJS,
PostgreSQL, vLLM, MCP, Stripe.

\end{document}
